% theory/definitions.tex -- a standalone LaTeX fragment. NOT wired into paper/main.tex.
%
% It assumes the preamble of paper/main.tex (amsmath, amssymb, amsthm and the
% environments definition, remark, proposition, theorem, proof) and uses two
% citation keys, troyanov1991 and thurston1998, that are in refs/sources.bib
% but not yet in the thebibliography of main.tex.
% Source records for those two keys: theory/teichmuller-dimension-sources.md.
% Numerical statements below are checked by theory/definitions-check.py.

\newcommand{\Pill}{\mathcal{P}}          % a comparison class of pillows
\newcommand{\Orb}{\mathcal{O}}           % a hyperbolic 2-orbifold
\newcommand{\Kiso}{K_{\mathrm{iso}}}
\newcommand{\Kmult}{K_{\mathrm{mult}}}

\subsection{Signatures and leading heat coefficients}\label{subsec:defs}

\begin{definition}[Signature]\label{def:signature}
Let $\Orb$ be a closed, orientable $2$-orbifold whose only singular points are cone points.
Let $g$ be the genus of the underlying surface and $m_1,\dots,m_n\ge2$ the orders of the cone
points, listed with multiplicity. The \emph{signature} of $\Orb$ is
\[
  \sigma(\Orb)=(g;\,m_1,\dots,m_n),
\]
in which the $m_i$ form a multiset: the order of listing carries no information. Write
$\chi(\Orb)=2-2g-\sum_{i=1}^n(1-\tfrac1{m_i})$. The orbifold carries a hyperbolic metric exactly
when $\chi(\Orb)<0$. A \emph{hyperbolic triangular pillow} is the case $(g;n)=(0;3)$; it is
denoted $\Orb(p,q,r)$ with $p\le q\le r$, and $\chi(\Orb(p,q,r))=R-1$ for
$R=\tfrac1p+\tfrac1q+\tfrac1r$. Two orbifolds are \emph{isometric} if there is a Riemannian
isometry of the hyperbolic metrics; isometric orbifolds have equal signatures.
\end{definition}

\begin{definition}[Leading heat coefficients]\label{def:heatcoef}
For a closed orientable hyperbolic $2$-orbifold $\Orb$ with Laplacian $\Delta$ the heat trace
has an asymptotic expansion in integer powers of $t$ \cite{donnelly1976, dggw2008},
\[
  \operatorname{tr}e^{-t\Delta}\;\sim\;\sum_{j\ge1}c_j(\Orb)\,t^{\,j-2}\qquad(t\downarrow0).
\]
The coefficient $c_j(\Orb)$ is the \emph{$j$-th leading heat coefficient}. In particular
$c_1(\Orb)=\operatorname{Area}(\Orb)/4\pi=-\chi(\Orb)/2$, which for a triangular pillow equals
$\tfrac12(1-R)$. For $k\ge1$ put $H_k(\Orb)=(c_1(\Orb),\dots,c_k(\Orb))$.
\end{definition}

\subsection{Two notions of how many coefficients suffice}

Fix a class $\Pill$ of closed orientable hyperbolic $2$-orbifolds. Everything below is relative
to $\Pill$; the statements of the manuscript concern $\Pill=\Pill_3$, the class of all hyperbolic
triangular pillows, and the class has to be named, because the answer changes with it.

\begin{definition}[$\Kiso$ and $\Kmult$]\label{def:K}
For $\Orb\in\Pill$ let
\begin{align*}
  \Kiso(\Orb;\Pill)&=\min\{k\ge1:\ \text{for every }\Orb'\in\Pill,\ H_k(\Orb')=H_k(\Orb)\Rightarrow\Orb'\ \text{is isometric to }\Orb\},\\
  \Kmult(\Orb;\Pill)&=\min\{k\ge1:\ \text{for every }\Orb'\in\Pill,\ H_k(\Orb')=H_k(\Orb)\Rightarrow\sigma(\Orb')=\sigma(\Orb)\},
\end{align*}
with $\min\emptyset=\infty$. Thus $\Kiso$ is the least number of leading heat coefficients that
determine $\Orb$ up to isometry among the members of $\Pill$, and $\Kmult$ is the least number that
determine its signature, i.e.\ its genus and its cone-order multiset.
\end{definition}

Both sets are upward closed in $k$, so the minima are well defined. Because isometric orbifolds
have equal signatures, the condition defining $\Kiso$ implies the one defining $\Kmult$, and
therefore
\begin{equation}\label{eq:Kcompare}
  \Kmult(\Orb;\Pill)\ \le\ \Kiso(\Orb;\Pill).
\end{equation}

\subsection{Rigidity: the two coincide for triangular pillows}

\begin{proposition}[Rigidity of triangular pillows]\label{prop:rigidity}
If $\Orb,\Orb'\in\Pill_3$ have the same signature then they are isometric. Consequently
\[
  \Kiso(\Orb;\Pill_3)=\Kmult(\Orb;\Pill_3)\qquad\text{for every }\Orb\in\Pill_3 .
\]
\end{proposition}

\begin{proof}
Label the cone points $x_1,x_2,x_3$ of $\Orb$ and $x_1',x_2',x_3'$ of $\Orb'$ so that $x_i$ and
$x_i'$ have the same order $m_i$. Let $\varphi$ be the M\"obius transformation of the sphere with
$\varphi(x_i)=x_i'$; it exists because the M\"obius group acts triply transitively. The pull-back
$\varphi^*g'$ of the hyperbolic cone metric $g'$ of $\Orb'$ is a conformal metric of curvature $-1$
with cone angle $2\pi/m_i$ at $x_i$. By Troyanov's Theorem~A \cite{troyanov1991}, applied with the
negative function $-1$ and the angles $\theta_i=2\pi/m_i$ (the hypothesis
$2\pi\chi(S^2)+\sum_i(\theta_i-2\pi)<0$ is exactly $\sum_i(1-1/m_i)>2$, i.e.\ hyperbolicity), such a
metric is unique in its conformal class. Hence $\varphi^*g'=g$ and $\varphi$ is an isometry.
Equal signature implies isometry, and isometry implies equal signature, so the two defining
conditions in Definition~\ref{def:K} agree, which gives the equality of $\Kiso$ and $\Kmult$.
\end{proof}

The reason is dimensional: the moduli space of hyperbolic cone metrics on the sphere with prescribed
cone angles is a point exactly when the sphere has three marked points, and has positive dimension
otherwise (next subsection).

\subsection{Positive-dimensional moduli force \texorpdfstring{$\Kiso=\infty$}{K-iso = infinity}}

By Troyanov's Theorem~A a hyperbolic cone metric on the sphere with prescribed angles is the same
datum as a conformal structure on the sphere with $n$ labelled marked points, and the space of such
structures has complex dimension $n-3$; Thurston's description of the space of shapes of
polyhedra with $n$ cone points, locally isometric to $\mathbb{C}H^{\,n-3}$, is the flat analogue
\cite{thurston1998}. Hence, for genus $0$,
\begin{equation}\label{eq:moduli}
  \dim_{\mathbb R}\mathcal M(0;m_1,\dots,m_n)=2n-6,
\end{equation}
which is $0$ for $n=3$ and positive for every $n\ge4$. (This is a two-step deduction from the two
cited sources and is flagged as such in \texttt{theory/teichmuller-dimension-sources.md}; for $g\ge1$
the corresponding count is not used here and would need its own source.)

The second ingredient is the following statement, which is \emph{not proved here}; it is the
subject of a separate locality theorem.

\begin{theorem}[Locality; forward reference]\label{thm:locality}
For every $j\ge1$ there is a function $\Phi_j$ of the signature alone with
$c_j(\Orb)=\Phi_j(\sigma(\Orb))$ for every closed orientable hyperbolic $2$-orbifold $\Orb$.
\end{theorem}

\noindent
Informally: heat invariants cannot see moduli. What is available now is the cone part: the
closed forms of Dryden--Gordon--Greenwald--Webb, Schueth and U\c{c}ar
\cite{dggw2008, schueth2019, ucar2017} depend only on the order $m_i$ of the cone point, for every
$j$ in U\c{c}ar's case, and were recomputed in exact arithmetic for $j\le6$
(\texttt{theory/cone-coefficients/}). What the locality theorem has to supply is the smooth part
at every order and the assembly of the pieces.

\begin{proposition}[Conditional on Theorem~\ref{thm:locality}]\label{prop:Kinf}
Suppose $\Pill$ contains two non-isometric orbifolds of the same signature $\sigma_0$. Then
$\Kiso(\Orb;\Pill)=\infty$ for every $\Orb\in\Pill$ with $\sigma(\Orb)=\sigma_0$. In particular,
if $\Pill\supseteq\mathcal M(0;m_1,\dots,m_n)$ with $n\ge4$, then $\Kiso(\Orb;\Pill)=\infty$ for all
$\Orb$ of signature $(0;m_1,\dots,m_n)$.
\end{proposition}

\begin{proof}
By Theorem~\ref{thm:locality} the two orbifolds have equal $H_k$ for every $k$, so no $k$ satisfies the
condition defining $\Kiso$. The second assertion follows from \eqref{eq:moduli}.
\end{proof}

\noindent
So $\Kiso$ is finite only where signature classes are single isometry classes, which among spheres
is the case $n=3$ and by Proposition~\ref{prop:rigidity} no other.
By \eqref{eq:Kcompare} this infinity is a property of $\Kiso$ alone: $\Kmult$ is unaffected.

\subsection{Theorem C and the larger-cone-count remark in this language}

\begin{theorem}[Theorem~\ref{thmC}, restated]\label{thm:Crestated}
For every $\Orb\in\Pill_3$,
\[
  \Kiso(\Orb;\Pill_3)=\Kmult(\Orb;\Pill_3)\le3 .
\]
The value is $3$ exactly when $\Orb$ belongs to a two-coefficient collision, that is, when some
$\Orb'\in\Pill_3$ with $\sigma(\Orb')\ne\sigma(\Orb)$ has the same $R$ and the same $S_1$; for
instance for $\Orb(2,8,8)$ and $\Orb(3,3,12)$.
\end{theorem}

\begin{proof}
The first three coefficients determine $(R,S_1,P_3)$ and hence the multiset $\{p,q,r\}$
(Proposition~\ref{prop:recovery}, the Newton--Vieta recovery), so $\Kmult\le3$; equality with
$\Kiso$ is Proposition~\ref{prop:rigidity}. The coefficients $c_1,c_2$ are equivalent to
$(R,S_1)$, so $\Kmult\le2$ unless a pillow of different signature shares $(R,S_1)$.
\end{proof}

\begin{remark}[Larger cone counts, restated]\label{rem:nconerestated}
Let $\Pill_n$ be the class of hyperbolic spheres with $n\ge3$ cone points, $\sum_i(1-1/m_i)>2$.
\begin{enumerate}[label=(\alph*)]
\item \emph{Multiset level.} Modulo the universal smooth contributions, $c_1,\dots,c_n$ are the $n$
  functions $R,\,S_1,\,P_3,\,P_5,\dots,P_{2n-3}$ of the multiset $\{m_i\}$, with $P_k=\sum_im_i^k$ and
  $c_j$ carrying $P_{2j-3}$ for $j\ge2$ with a nonzero weight. If the map
  $\{m_i\}\mapsto(R,S_1,P_3,\dots,P_{2n-3})$ is injective on $n$-element multisets, then
  $\Kmult(\Orb;\Pill_n)\le n$. For $n=3$ this is Proposition~\ref{prop:recovery}; for $n\ge4$ injectivity is not
  proved here.
\item \emph{Isometry level.} For $n\ge4$, $\Kiso(\Orb;\Pill_n)=\infty$ by
  Proposition~\ref{prop:Kinf} and \eqref{eq:moduli}. The statement that the bound $K\le n$ extends
  the case $n=3$ of Theorem~\ref{thmC} is therefore valid for $\Kmult$ (conditionally on injectivity) and
  false for $\Kiso$; what singles out $n=3$ is rigidity, Proposition~\ref{prop:rigidity}.
\item \emph{A lower bound at $n=4$.} The multisets $\{3,10,15,30\}$ and $\{4,5,21,28\}$ are
  hyperbolic, have the same $S_1=58$, $R=8/15$ and $P_3=31402$, and have $P_5=25159618\ne21298618$
  (exact computation, \texttt{theory/definitions-check.py}). Hence $\Kmult\ge4$ for pillows with
  these signatures, and $\Kmult=4$ for them if the injectivity in (a) holds for $n=4$.
\end{enumerate}
\end{remark}

% Nonvanishing of the weight in (a) for every j rests on the Bernoulli-number closed form for the
% leading coefficient of p_l; it is verified exactly for l = 0..6 in
% theory/cone-coefficients/verify_cone_coefficients.py and is not proved for all l here.
